% !TeX root = ../main.tex

\begin{abstract}[zh]
  碱性电解水（AWE）是大规模绿氢制备的主流技术，但风电、光伏等波动性电源的并网
  对多槽并联AWE系统的协调控制提出了更高要求。传统非线性模型预测控制（NMPC）
  计算开销大，难以满足实时性需求；而纯数据驱动的学习方法又缺乏对温度、氢气纯度
  等安全约束的严格保证。本文面向四槽并联AWE系统的功率分配与温度调节问题，
  建立了一套基于扩散模型的安全预测控制框架。

  研究首先建立了涵盖电化学、热力学与氢气传输的高保真仿真模型。
  在此基础上，以NMPC为专家控制器生成训练数据并训练TCN Diffusion控制策略，
  进而引入控制障碍函数（CBF）安全滤波机制，
  通过CBF理论将生成策略的无约束输出修正为满足安全边界的可行动作，
  有效抑制扩散模型推理过程中偶发异常输出带来的安全隐患。实验结果表明，
  所提策略的功率跟踪精度接近NMPC基线，推理时间从NMPC的约17.6~s降至约177~ms，
  实现了近两个数量级的计算加速，且HTO始终低于2\%的安全阈值，
  温度、电压与功率约束满足率均达到100\%。

\end{abstract}

\begin{abstract}[en]
  Alkaline Water Electrolysis (AWE) is the mainstream technology for large-scale
  green hydrogen production. However, the integration of variable renewable energy
  sources such as wind and solar power poses higher demands on the
  coordinated control of multi-stack parallel AWE systems. Traditional Nonlinear
  Model Predictive Control (NMPC) suffers from heavy computational overhead and
  fails to meet real-time requirements, while purely data-driven approaches lack
  strict guarantees for safety constraints such as temperature and hydrogen purity.
  To address these issues, this study focuses on four-stack parallel AWE systems and
  proposes a safe predictive control framework based on diffusion models.

  First, we establish a high-fidelity simulation model covering electrochemistry,
  thermodynamics, and hydrogen transport. Then, we employ NMPC as an expert
  controller to generate training data and train a TCN Diffusion control
  policy on the resulting trajectories.
  Finally, we introduce Control Barrier Functions (CBF) as a safety filter,
  leveraging Lie-derivative forward invariance theory to correct the unconstrained
  policy outputs into safe feasible actions, thereby effectively suppressing the
  tail risk inherent in diffusion model inference. Experimental results show that the proposed strategy achieves
  power tracking accuracy close to the NMPC baseline while reducing inference
  time from about 17.6~s to about 177~ms, yielding a computational speedup of nearly two orders of magnitude. HTO remains strictly below the 2\% safety
  threshold, and the temperature, voltage, and power constraints are fully satisfied.

\end{abstract}
